\documentclass[12pt]{article}
\usepackage[utf8]{inputenc}
\usepackage[T1]{fontenc}
\usepackage[margin=1in]{geometry}
\usepackage{setspace}
\usepackage{xcolor}
\usepackage[colorlinks=true, urlcolor=blue!60!black, linkcolor=blue!60!black]{hyperref}
\setlength{\parindent}{0em}
\setlength{\parskip}{0.45em}
\renewcommand{\baselinestretch}{1.0}
\begin{document}
\begin{center}
{\Large \textit{Seila Law LLC v. CFPB}, 591 U.S. 197 (2020)\par}
\vspace{0.3em}
{\normalsize Case Brief --- One-Page Edition\par}
\end{center}

\textbf{Memory Jogger.} A single agency head wielding significant executive power must be removable at will; the CFPB's for-cause protection was unconstitutional but severable (\href{https://www.law.cornell.edu/supremecourt/text/19-7}{Cornell LII}).

\textbf{Facts.} Dodd-Frank created the CFPB with broad rulemaking, enforcement, and adjudicatory power under a single Director with a five-year term, removable only for ``inefficiency, neglect of duty, or malfeasance in office,'' and funded through the Federal Reserve. Seila Law, a debt-relief law firm, resisted a civil investigative demand on the ground that this structure was unconstitutional.

\textbf{Procedural History.} The district court enforced the demand; the Ninth Circuit affirmed under \textit{Humphrey's Executor}.

\textbf{Votes.} Vacated and remanded: 5--4 on removal (Roberts; Kagan dissenting for four) and 7--2 on severability (Thomas and Gorsuch dissenting).

\textbf{Holding.} Executive power generally includes at-will removal, with exceptions only for multimember expert bodies and limited inferior officers. The Director's protection is invalid but severable.

\textbf{Reasoning.} The CFPB fits neither exception; a single unaccountable Director has little historical precedent and defies a design that concentrates accountability in the elected President, aggravated by the long term and independent funding. Dodd-Frank's severability clause keeps the agency running. Thomas would overrule \textit{Humphrey's Executor} and deny enforcement outright; Kagan answered that the Necessary and Proper Clause lets Congress insulate regulators and that a single Director is easier to oversee than a board.

\textbf{Examples.} \emph{Same-Side:} the single FHFA Director (\textit{Collins v. Yellen}). \emph{Opposite-Side:} a bipartisan adjudicatory commission. \emph{Fence-Sitter:} the modern FTC, later addressed in \textit{Trump v. Slaughter} (2026).

\textbf{Critique.} Revisionist historians dispute the Court's account of the Decision of 1789, and functionalists call the single-head line arbitrary.

\textbf{Key Quotation.} ``Under our Constitution, the `executive Power'---all of it---is `vested in a President,' who must `take Care that the Laws be faithfully executed.' Art. II, \S 1, cl. 1; id., \S 3.'' --- Chief Justice Roberts, for the Court.

\end{document}
